% ===== §6 =====
\section{The Generativity That Justice Guards}
\label{sec:justice}

\noindent
This section takes the fifth dimension, the justice a relation owes, and establishes that what justice most owes protection to is not only a settled fact awaiting adjudication but the generativity of relational being, the still-forming life of the bond, which is the intervention of the Real into the symbolic and cannot be settled in advance by any existing order of signs. The result is that justice at its root is the guarding of generativity rather than the settlement of the already-formed. The section proceeds through the phenomenon, a review of the two pictures of justice that treat it as settlement, the isolation of the core, and the claim.

\subsection*{The phenomenon}

A relation owes justice, and justice seems the most sayable of all the dimensions, since it lives in claims, procedures, and judgments, in the very element of statement. Justice weighs what is owed and settles it, and to settle is to say. It appears that whatever justice concerns can in principle be brought into the order of signs, stated as a claim and disposed of by a judgment. Two pictures reinforce the appearance, and it is worth naming them to see what they leave out.

\subsection*{Two pictures of justice as settlement}

The first is the procedural picture, on which justice is the correct running of a procedure that takes claims as inputs and yields judgments as outputs. Its virtue is that it makes justice sayable through and through: everything relevant is a claim, and a claim is a statement. But a procedure operates on the already-formed. It takes what has happened, states it as a case, and disposes of it, and it is constitutively backward-looking, addressed to what has taken shape and can be brought before it. What is still forming, what has not yet taken the shape a claim would state, is not an input the procedure can accept, and so it falls outside the procedure's reach not by oversight but by the procedure's nature.

The second is the picture of justice as the retrospective repair of wrongs, the setting-right of a harm already done. This too is real and necessary, and this too is backward-looking. It addresses the wrong as a fact, laid out for assessment, and its whole labor is the disposal of what has already occurred. Here the work on testimony is instructive by contrast. Felman and Laub, studying the witnessing of catastrophe, found that the most severe wrongs resist exactly the settlement that repair would require, that there is in them a core which does not become a statable fact to be assessed and closed, but persists as something only borne witness to, never disposed of \citep{felman1992testimony}. Even in the retrospective case, then, a remainder escapes settlement. And what escapes there is a version of what concerns us, though the section's own subject lies not in the past wrong but in the forward direction, in the still-forming.

\subsection*{The core}

Something in justice resists settlement, and it is the deepest thing justice has to protect. What a relation of justice must answer to is not only a settled fact but the generativity of relational being itself, the still-forming, not-yet-settled life of the bond. This generativity is the intervention of the Real into the symbolic. What is new is new precisely in that it is not yet covered by the existing signifiers, a tear the Real opens in the order of established meaning. The core language cannot reach here is this \emph{generativity} of relational being, the Real's intervention, which cannot be settled in advance by any existing order of signs. To bring the still-forming under the settled sign is to demote it from becoming to become, from the living generativity that justice must protect to a finished case that a procedure can dispose of. The limit is not that justice has not yet found the right words. It is that the object justice most owes protection to is, by its nature as generativity, what the words have not reached and cannot reach without ending it.

\begin{claim}[Justice guards what it cannot settle]
The deepest object of justice is the generativity of relational being, the still-forming life that the intervention of the Real keeps opening in the order of settled meaning. This generativity cannot be settled in advance, for to settle it is to convert becoming into the already-become and so to end it. Justice at its root is therefore the guarding of a generativity it cannot bring into the settled sign, not the settlement of the already-formed alone.
\end{claim}

\begin{casebox}{the grievance that the settlement misses}
A wrong within a bond is brought to a reckoning, stated as a claim, weighed, and disposed of by an agreed repair. The repair is fair, and something in the wrong is nonetheless left untouched by it, a residue that the statement of the claim could not carry and the settlement could not reach. What the reckoning addressed was the wrong insofar as it could be stated; what remains is the wrong insofar as it could not. The residue is not a sign that the settlement was unjust. It is the mark that justice answers to something a settlement does not exhaust.
\end{casebox}

\subsection*{Why the limit is generative here}

This is what keeps justice from collapsing into administration. Because generativity cannot be settled in advance, justice cannot be exhausted by procedure, and a justice worthy of the name holds itself open to what its own settled forms have not anticipated. Consider the counter-case, a justice that could dispose of everything by existing rule. It would be an administration of the already-formed, complete and closed, and precisely in its completeness it would be deaf to the new that relational being keeps generating, since the new is what no existing rule anticipated. The limit forbids this closure. Because the generativity justice guards cannot be brought under the settled sign, justice must keep itself answerable to what exceeds its procedures, and refuse itself the false completeness of a procedure that has foreseen everything. The mode of this keeping-open, a knowing that holds what it cannot dispose of, is taken up in Part Four under the name of witnessing. What justice cannot bring into the settled sign is not a gap in the law. It is the living generativity that the law exists to keep open.
